\subsection{Figures, and the results that motivate the rest}
\label{sec:figures}

The parameter values are illustrative; none is fitted. They establish qualitative
structure and orders of magnitude, not any particular patient, and
Table~\ref{tab:params_orig} says which of them are derived and which are assumed.

The ordering of this section is deliberate and it changed in the restructure. The
fibre, the factor of $169$ and the shortfall of the accessible perturbations used to be
presented as the conclusions of the paper. They are not: they are corollaries of the
local reduction, and what they establish is that a static reading of the record cannot
answer the question, which is why the dynamic model and the excitation condition exist.
Nothing has been deleted. What has changed is what the figures are for.

%-------------------------------------------------------------------
\subsubsection{Corollaries of the local reduction: the degeneracy and what it costs}

The two elimination channels, and the point-matched pair that no affine fit separates,
are Figure~\ref{fig:A1} of the paper.

\begin{figure}[htbp]
\centering
\begin{minipage}[t]{0.48\textwidth}
\centering
\includegraphics[width=\textwidth]{fig_A2_conditioning.pdf}
\end{minipage}\hfill
\begin{minipage}[t]{0.48\textwidth}
\centering
\includegraphics[width=\textwidth]{fig_A5_load_error.pdf}
\end{minipage}
\caption{Left: condition number of the column-normalised sensitivity matrix against
the relative excursion of the metabolite, on log axes. The slope is $-2$: conditioning
degrades quadratically as the patient is stabilised, and a patient varying by a few
per cent between visits presents a problem conditioned at $10^{5}$ to $10^{6}$.
Right: an error in the assumed substrate load propagates into $\theta$ amplified by
$C(\bar m)/\theta g(\bar m)$, total clearance over enzymatic clearance, because the
enzymatic term is recovered as a difference of two larger quantities. That ratio is not a
property of the model. It is $1.4$ at these illustrative values only because most of the
flux goes through the enzyme, it reaches $7$ at $\theta = 100$, and it diverges as
$\theta \to 0$, so the amplification is worst for exactly the patients the paper is
about.}
\label{fig:A2}
\end{figure}

\begin{figure}[htbp]
\centering
\includegraphics[width=\textwidth]{fig_A3_indistinguishable.pdf}
\caption{Left: a selective perturbation renders two patients matched at the operating
point formally distinguishable. The perturbation shown is carnitine, which acts on the
alternative channel $k_e$ and is administered to every patient in the accessible cohort.
The separation remains near $1\%$ across the plausible dose range, well below assay
variability: formal identifiability and practical identifiability part company here.
Right: the separation plotted against the excursion of $m$ that the perturbation actually
achieves, for both channels. The two curves lie on top of one another, which is
Proposition~\ref{prop:channel} made visible. The excursion is bought by moving the patient
and it does not matter which channel is used to move them, which is what allowed the
analysis to be posed on carnitine while the therapy was prospective. Two things have since
changed and the caption is read with them. The therapy is administered, so the analysis is
posed on it directly (Section~\ref{sec:therapychannel}). And the equality holds at fixed
$K_m$: with $K_m$ free the two routes separate by six orders of magnitude in $\det G$,
because only the therapy modulates the prefactor that carries the $K_m$ sensitivity
(Remark~\ref{rem:channel_scope}). What this figure shows is the excursion, not the
curvature.}
\label{fig:A3}
\end{figure}

\begin{figure}[htbp]
\centering
\includegraphics[width=\textwidth]{fig_A4_population.pdf}
\caption{Left: \textbf{one steady state per patient}. A cohort of $N$ patients sharing the
population parameters yields $N$ equations in $N+2$ unknowns, and the deficiency is two for
every $N$, so enlarging the cohort adds an equation and an unknown together. This is the
counting for a single observation each, and it is not the data model
Proposition~\ref{prop:degeneracy} is built on. A patient observed over a band supplies two
coefficients, giving $2N$ equations in $N+2$ unknowns, determined at $N = 2$ and
overdetermined beyond: two patients alone give
$K_m = (b_2 - b_1)/(a_2/\bar m_2^2 - a_1/\bar m_1^2)$, whenever the denominator is
non-zero, without either patient moving. What limits cohort closure is therefore the
conditioning of the intercepts and not the counting. Right: with two anchors the
single-observation system closes exactly. A null genotype supplies $k_e$ directly, since
its equation loses the enzymatic term; one patient observed at two distinct loads then
supplies $K_m$; every remaining $\theta_i$ follows from a single equation. Recovery is
exact in the noise-free case.}
\label{fig:A4}
\end{figure}

\begin{figure}[htbp]
\centering
\includegraphics[width=\textwidth]{fig_A6_fibre.pdf}
\caption{The fibre a perfectly resolved band still leaves open,
Remark~\ref{rem:fibre}. Left: the residual capacities consistent with a known $(a,b)$,
against the $K_m$ that indexes them. With noiseless data $\theta$ spans a factor of $169$,
and the true patient is one point on that curve. Right: the same fibre in $k_e$, which
reaches zero at $K_m = b\bar m^2/a$ and bounds the family from above. A null genotype fixes
$k_e$ and collapses the fibre to a point, recovering $K_m$ and $\theta$ exactly and without
a second steady state. The band that resolves $(a,b)$ does the work that
Proposition~\ref{prop:anchor}'s second anchor was introduced for.}
\label{fig:A6}
\end{figure}

\begin{figure}[htbp]
\centering
\includegraphics[width=\textwidth]{fig_A7_margin.pdf}
\caption{The decision quantity is better conditioned than the parameter it is built from,
and where that stops being enough. (a) The safety margin
$\mu = \theta - \thc^{\mathrm{A}}$ along the same fibre. Where $\theta$ is quadratic in
$K_m$ the margin is linear, the quadratic terms cancelling identically in
\eqref{eq:mufibre}, so the whole sensitivity to the unresolved direction is one slope,
computable from the fitted $(a,b)$. The improvement is one order, and $\mu$ still varies
several-fold across the shaded range. (b) Because $\mu$ is linear its
sign is constant along the fibre unless the two endpoints disagree, which happens only for
the shaded interval of illness amplitudes, \eqref{eq:window}. Outside it the decision is
determined although $\theta$ is not; inside it no width of band settles the decision.
Illness amplitude is drawn as a multiple of the patient's own baseline load, and the dotted
line is the amplitude the equilibrium reduction applies in Figures~\ref{fig:C1} and
\ref{fig:C2}.
(c) The width of that window, against the convention that sets it. It closes as
$\mtol$ approaches the band the patient is observed in and widens without bound above it.
(d) Two patients matched at the operating point, whose margins separate exactly as
$\Delta k_e\,(\mtol - \bar m)$ and coincide only when the tolerable level is the operating
point itself. The same convention that widens the window separates the patients.}
\label{fig:A7}
\end{figure}

\begin{figure}[htbp]
\centering
\includegraphics[width=\textwidth]{fig_A8_design.pdf}
\caption{What the design would have to deliver, against what it does, all four panels
computed from \eqref{eq:relerr} and its companions in Section~\ref{sec:design}.
(a) Equation \eqref{eq:thetaerr}, evaluated, as a family over how well the
substrate load is recorded, since nothing in the paper bounds that. The shaded region is what
the quasi-exogenous instruments of Section~\ref{sec:instruments} reach, and every curve
crosses the $20\%$ target well outside it. (b) The band at which the curvature of
$C$ becomes resolvable at two standard errors, which is when condition \eqref{eq:affine} can
be checked against a record, against the same load error. It too lands
on the far side of the same line, from about $2.5$ times outside at a load known to $2\%$ to
about $8$ times at $20\%$. (c) The steady-state design puts the noise in the
abscissa. An unconstrained fit appears to flip the sign of the margin, but it does so by
returning $\hat k_e < 0$, which is not a parameter set; with $k_e \ge 0$ imposed the bias is
small. The grey curve is the fraction of unconstrained fits that are inadmissible, and it is
what makes the failure detectable from the fit alone. Remark~\ref{rem:eiv}. (d) The
cost of the null anchor is a ratio, \eqref{eq:anchorfac}, not a constant. It attenuates at
$\theta = 3000$ and amplifies below $\theta = 1100$, so it is worst for the severe patients
the paper is about, and the parent's own $\theta$ sits on the amplifying side.}
\label{fig:A8}
\end{figure}

\begin{proposition}[Anchor condition]
\label{prop:anchor}
Let $N$ patients at steady state share $(K_m, k_e)$. If the cohort contains at least
one patient with a null genotype, so that $\theta_j = 0$, and at least one patient
observed at two distinct steady states, then $K_m$, $k_e$ and every $\theta_i$ are
uniquely determined.
\end{proposition}

Both anchors are present in the Saudi series: the frameshift variant \emph{PCCB}
c.990dupT occurs in seven families and is a null allele~\cite{saudi84}, and documented changes in % C26
prescribed protein intake supply the second condition.

The second condition can be dispensed with. Proposition~\ref{prop:anchor} counts one
observation per patient. Under the band model of Proposition~\ref{prop:degeneracy} each
patient supplies two coefficients, and by Remark~\ref{rem:fibre} a null genotype's $k_e$
then closes the fibre without anyone being moved to a second steady state. Two patients with
bands and no null genotype suffice as well, since $(a_i, b_i)$ for $i = 1, 2$ give $K_m$
directly. What the excursion buys is not the counting but the conditioning of the
intercepts, which is the subject of \eqref{eq:thetaerr}.

What justifies setting $\theta_j = 0$ for the null patients is the form of the model, and needs no
claim about residual enzyme activity. The construction
does not require that \emph{no} propionate leaves by any other route. It requires only
that whatever else eliminates propionate is not $\theta\, g(m)$. Any residual elimination,
of whatever molecular origin, enters \eqref{eq:A} as part of the alternative channel and is
absorbed into $k_e$, which is precisely the quantity the null patients measure. % C27
In particular the anchor does not rest on any claim about the molecular origin of residual
carboxylation in PCC-deficient cells. Whether it comes from a distinct carboxylase or from
partially functional mutant enzyme is a question this construction never has to answer.

\clearpage
%-------------------------------------------------------------------
\subsubsection{The alternative hypothesis with a ceiling, and what it costs}

\begin{figure}[htbp]
\centering
\begin{minipage}[t]{0.49\textwidth}
\centering
\includegraphics[width=\textwidth]{fig_B1_ceiling.pdf}
\end{minipage}\hfill
\begin{minipage}[t]{0.49\textwidth}
\centering
\includegraphics[width=\textwidth]{fig_B2_bifurcation.pdf}
\end{minipage}
\caption{Left: with both channels saturating, total clearance is bounded by
$\theta + k_e$. Loads below the ceiling admit exactly one steady state; loads above it
admit none. Right: the equilibrium as a function of load. It escapes to infinity as the
load approaches the ceiling rather than colliding with a second branch, so this is not
a saddle-node. The clinically operative threshold arrives earlier, when the equilibrium
migrates past the tolerable level.}
\label{fig:B1}
\end{figure}

\begin{figure}[htbp]
\centering
\begin{minipage}[t]{0.47\textwidth}
\centering
\includegraphics[width=\textwidth]{fig_B3_thetacrit.pdf}
\end{minipage}\hfill
\begin{minipage}[t]{0.50\textwidth}
\centering
\includegraphics[width=\textwidth]{fig_B4_cost.pdf}
\end{minipage}
\caption{Left: $\thc$ over the plane of the two quantities it depends on, neither of
which is measured. The critical capacity is an operational definition; no physiology fixes
it, and the figure shows how much of its value is inherited from
that choice. Right: the price of biological fidelity, of which this figure shows only the
local half. Once both channels saturate they share a functional form and are separated
only by differing half-saturation constants; at a ten per cent excursion the condition
number rises from $1.1 \times 10^{5}$ in the local reduction to $4.8 \times 10^{7}$ under
the ceiling hypothesis, a factor of roughly $450$. A condition number cannot show the
other half: by Proposition~\ref{prop:exchange} the ceiling hypothesis also carries an
\emph{exact} two-point degeneracy in which $(\theta,K_m)$ and $(k_e,K_e)$ swap. The
Jacobian stays full rank, which is why a finite condition number is the right local
description, and a second isolated parameter vector reproduces the data exactly all the
same.}
\label{fig:B2}
\end{figure}

\begin{remark}
This is the central trade-off of the modelling exercise: the more faithful model is
the harder one to identify. The local reduction is retained as the working model for the
identifiability analysis because it represents the most favourable case: if identification
fails there it fails a fortiori once both channels bend.
\end{remark}

\clearpage
%-------------------------------------------------------------------
\subsubsection{The equilibrium reduction: depletion, a testable prediction, and a threshold}
\label{sec:resultsC}

The figures below are computed on the reduced clearance
$C(m;\theta) = \theta g(m) + \alpha(m)$ of \eqref{eq:alpha}, obtained from the two-state
system \eqref{eq:C} by solving $\dot c = 0$. What makes $\alpha$ vanish with the metabolite
and then decline is the second exit for carnitine and the widening of it,
Section~\ref{sec:model_C_added} setting out why stoichiometric consumption on its own would
return a ceiling instead.

\begin{figure}[htbp]
\centering
\begin{minipage}[t]{0.49\textwidth}
\centering
\includegraphics[width=\textwidth]{fig_C1_phaseline.pdf}
\end{minipage}\hfill
\begin{minipage}[t]{0.49\textwidth}
\centering
\includegraphics[width=\textwidth]{fig_C2_fold.pdf}
\end{minipage}
\caption{Left: the phase line under the equilibrium reduction. At baseline load the system has two
equilibria, a stable attractor at low metabolite and an unstable one at high
metabolite which acts as a separatrix. Under peak load neither exists. Right: the
corresponding bifurcation diagram. The two branches approach and collide, which is a
genuine saddle-node, and a ceiling admits no such collision.}
\label{fig:C1}
\end{figure}

\begin{figure}[htbp]
\centering
\begin{minipage}[t]{0.48\textwidth}
\centering
\includegraphics[width=\textwidth]{fig_C3_timeseries.pdf}
\end{minipage}\hfill
\begin{minipage}[t]{0.50\textwidth}
\centering
\includegraphics[width=\textwidth]{fig_C4_basin.pdf}
\end{minipage}
\caption{Left: the discriminating prediction. An identical transient illness is applied
to both structures. Under a ceiling the metabolite returns to baseline once the illness
resolves. Under the fold the patient has crossed the separatrix and continues to
deteriorate at baseline load. Right: the basin of attraction. There is no upper stable
state, so the clinically meaningful object is the separatrix. The marker is a patient
above the separatrix at baseline load, and the two arrows
are the routes out. Arrow~1 drives the load below baseline, which raises the separatrix
above the patient, and is why suppression of catabolism by intravenous glucose is the
mainstay of treatment. Arrow~2 lowers the metabolite directly, by haemofiltration, and is
the alternative when the load cannot be reduced far enough. The panel carries pointers
only, four labelled annotations having been unreadable at this size.}
\label{fig:C2}
\end{figure}

\begin{remark}[The discriminant, and what it cannot be measured in]
The two structures differ in a prediction: under a ceiling, returning $s$ to baseline returns $m$ to
baseline, and under C it does not, because the separatrix at baseline load still lies below
the patient. That difference is a statement about the model's state $m$, and the obstacle
is that $m$ cannot be watched through a crisis. Deep phenotyping of acute decompensation
found propionylcarnitine, propionylglycine and propionic acid essentially unaltered during
the event~\cite{amd_phenotyping}. % C24
The single marker does not track the state through the excursion the discriminant is about,
so ``does 3-HP return to baseline'' is not the question to take to the records. Free
carnitine does move, its median falling from $38.9$ to $22.0$ $\mu$M between samples
drawn outside a decompensation and samples drawn during one, $P =
0.0370$~\cite{amd_phenotyping}, % C38
and Section~\ref{sec:discriminant} restates the discriminant in it, carrying with it the
group-level character of that comparison and the fact that C is not the only model it
fits.
\end{remark}

\begin{figure}[htbp]
\centering
\includegraphics[width=\textwidth]{fig_C5_margin.pdf}
\caption{The intrinsic critical capacity, Proposition~\ref{prop:intrinsic}, computed at
baseline load. Left: the attractor and the separatrix as functions of the residual
capacity $\theta$, on a logarithmic metabolite axis. As $\theta$ falls the attractor rises
and the separatrix falls, and at $\thc^{\,\mathrm{int}} = 151$ they merge in a saddle-node;
below that value no equilibrium exists at all and the patient cannot be maintained even at
baseline load. Nothing in this construction refers to a tolerable metabolite level: the
threshold is a property of the model's equilibrium structure, not of a clinical convention.
Right: the distance between the two, which is the safety margin. It is the quantity the
clinical decision actually turns on, since it says how far the patient can be pushed before
the crisis becomes self-sustaining, and it closes continuously to zero at the threshold.
It is drawn as a band and not as a curve, and that is the honest form of it. The
separatrix sits where $\alpha$ has decayed, and $\alpha$ decays on the scale of $M_i$,
which is a constant nobody has measured; the band is $M_i$ swept over an eight-fold range
with $R_0$ recalibrated at each point, so no admissible parameter set is broken to draw it.
The sign is robust and the magnitude is not: the margin stays positive at every
$\theta$ and every $M_i$ in the sweep, while at the nominal $\theta = 300$ it runs across
roughly a factor of six about the nominal $954$. Read the figure for the sign and for the
order of magnitude, and not for the value (Remark~\ref{rem:shape}). Neither reduced
structure admits this figure at all: A never loses its equilibrium, and B loses one to
infinity rather than to a collision, so in neither is there a second equilibrium to measure
a margin against.}
\label{fig:C5}
\end{figure}

\begin{remark}[The price of the operational threshold]
Figure~\ref{fig:B2}~(left) is the price of the operational $\thc$ of \eqref{eq:thcrit}, and
Figure~\ref{fig:C5} is what replaces it. How much of the operational threshold's numerical
value is inherited from the convention is visible in the first; the second is a location in
the model's own equilibrium structure, and Section~\ref{sec:intrinsic} sets out why only
only the equilibrium reduction admits it.
\end{remark}

